\documentclass[11pt]{article}

\usepackage[margin=1in]{geometry}
\usepackage{amsmath,amssymb}
\usepackage{booktabs}
\usepackage{graphicx}
\usepackage{hyperref}

\title{Econ 5253 -- Problem Set 12}
\author{Cheng}
\date{May 7, 2026}

\begin{document}
\maketitle

\section*{1--3.\ Repository setup and packages}
The required commands (\texttt{git pull origin master}, \texttt{git fetch upstream})
were issued from the OSCER \texttt{DScourseS26} folder, and the R packages
\texttt{sampleSelection}, \texttt{tidyverse}, and \texttt{modelsummary} were
installed.  All analysis is reproduced by the script
\texttt{PS12\_Cheng.R}.

\section*{4--6.\ Data summary and missingness}
The data set \texttt{wages12.csv} contains $N = 2{,}229$ observations on
working women in the U.S.\ in 1988.  After converting \texttt{college},
\texttt{married}, and \texttt{union} to factors, Table~\ref{tab:summary}
reports descriptive statistics.

\begin{table}[h]
    \centering
    \caption{Summary statistics for the \texttt{wages12} data}
    \label{tab:summary}
    \begin{tabular}{lrrrrrr}
        \toprule
        Variable     & Unique & Missing \% & Mean   & SD    & Min    & Max    \\
        \midrule
        \texttt{logwage} & 1546 & 30.69 &  1.652 & 0.688 & $-0.956$ &  4.166 \\
        \texttt{hgc}     &   14 &  0.00 & 12.455 & 2.444 &  5.000 & 18.000 \\
        \texttt{exper}   & 1932 &  0.00 &  6.435 & 4.867 &  0.000 & 25.000 \\
        \texttt{kids}    &    2 &  0.00 &  0.429 & 0.495 &  0.000 &  1.000 \\
        \midrule
        \multicolumn{7}{l}{\emph{Categorical variables (counts and percentages)}} \\
        \texttt{college} = Yes & \multicolumn{2}{r}{233 (10.5\%)} & \multicolumn{4}{l}{(No: 1996, 89.5\%)}  \\
        \texttt{married} = Yes & \multicolumn{2}{r}{1415 (63.5\%)} & \multicolumn{4}{l}{(No: 814, 36.5\%)}  \\
        \texttt{union}   = Yes & \multicolumn{2}{r}{529 (23.7\%)} & \multicolumn{4}{l}{(No: 1700, 76.3\%)} \\
        \bottomrule
    \end{tabular}
\end{table}

\paragraph{Discussion.}  The numbers look plausible for a 1988 cross-section
of working women: average completed schooling is about 12.5 years (roughly a
high-school graduate), average labour-market experience is 6.4 years, about
two-thirds of the sample is married, and 23.7\% are union members.  Mean log
wage is 1.65 (a wage of about \$5.20/hour, consistent with 1988 dollars).

\paragraph{Missingness.}  Log wages are missing for 684 of 2229 women, a
\textbf{30.7\%} missing rate.  Because the sample only contains women who
were working, the missingness in \texttt{logwage} is most likely
\textbf{Missing Not At Random (MNAR)}.  A woman with a low offered wage is
more likely to drop out of the labour force or refuse to report a wage,
i.e.\ her wage is missing precisely because of its value.  This pattern of
self-selection into reporting violates the MCAR and MAR assumptions and is
exactly the type of problem the Heckman correction was designed for.

\section*{7.\ Three imputation methods for the wage regression}
For each method I estimate
\[
\text{logwage}_i = \beta_0 + \beta_1 \text{hgc}_i + \beta_2 \text{union}_i +
                   \beta_3 \text{college}_i + \beta_4 \text{exper}_i +
                   \beta_5 \text{exper}_i^2 + \eta_i.
\]
The three columns of Table~\ref{tab:reg} correspond to (a) listwise deletion
of complete cases, (b) mean imputation of \texttt{logwage}, and (c) the
two-step Heckman selection model whose selection equation also includes
\texttt{married} and \texttt{kids} (the exclusion restrictions).

\begin{table}[h]
    \centering
    \caption{Returns to schooling under three missing-data treatments}
    \label{tab:reg}
    \begin{tabular}{lccc}
        \toprule
        & Listwise deletion & Mean imputation & Heckit (outcome) \\
        \midrule
        (Intercept)        &  0.834***  &  1.149***  &  0.446***  \\
                           & (0.113)    & (0.078)    & (0.122)    \\
        \texttt{hgc}       &  0.059***  &  0.036***  &  0.091***  \\
                           & (0.009)    & (0.006)    & (0.010)    \\
        \texttt{unionYes}  &  0.222*    &  0.068     &  0.186*    \\
                           & (0.087)    & (0.047)    & (0.084)    \\
        \texttt{collegeYes}& -0.065     & -0.126**   &  0.092     \\
                           & (0.106)    & (0.048)    & (0.100)    \\
        \texttt{exper}     &  0.050***  &  0.021**   &  0.054***  \\
                           & (0.013)    & (0.007)    & (0.012)    \\
        \texttt{exper}$^2$ & -0.004**   & -0.001**   & -0.002+    \\
                           & (0.001)    & (0.000)    & (0.001)    \\
        Inverse Mills      &            &            & -0.695***  \\
                           &            &            & (0.060)    \\
        \midrule
        $N$                & 1545     & 2229     & 2229 (684 cens.) \\
        $R^2$              & 0.038    & 0.020    & 0.092    \\
        \bottomrule
    \end{tabular}\\[2pt]
    \footnotesize{$+\,p<0.1$, $*\,p<0.05$, $**\,p<0.01$, $***\,p<0.001$.
    Standard errors in parentheses.  The Heckit selection equation also
    contains \texttt{married} and \texttt{kids}.  See \texttt{regression\_table.tex}
    for the raw \texttt{modelsummary} output.}
\end{table}

\paragraph{Comparison of $\hat{\beta}_1$.}
The true return-to-schooling parameter is $\beta_1 = 0.091$.  Across the
three estimators we obtain
\[
   \hat{\beta}_1^{\text{listwise}} = 0.0590, \qquad
   \hat{\beta}_1^{\text{mean}}     = 0.0363, \qquad
   \hat{\beta}_1^{\text{Heckit}}   = 0.0915.
\]
Both naive imputation methods produce a substantial \emph{downward} bias.
Listwise deletion under-estimates the return by about a third, and mean
imputation does even worse: replacing missing wages with the sample mean
mechanically shrinks the variance of the dependent variable around its mean
and pulls every slope toward zero, so $\hat{\beta}_1$ collapses to roughly
$0.036$.  In contrast the Heckman two-step procedure recovers
$\hat{\beta}_1 = 0.0915$, essentially identical to the true value of
$0.091$.  The strongly significant inverse-Mills coefficient
($-0.695$, $p < 0.001$) confirms that the missing-wage process is non-random
and that ignoring it produces biased OLS estimates.

\paragraph{Take-away.}  When data are MNAR, listwise deletion and mean
imputation are both unreliable; mean imputation is especially dangerous
because it injects fictitious observations at the centre of the
distribution.  Modelling the selection mechanism explicitly (Heckit) is the
only one of the three methods that delivers an unbiased estimate of the
returns to schooling here.

\section*{8.\ Probit model of union-job preferences}
I estimate
\[
   u_i = \alpha_0 + \alpha_1 \text{hgc}_i + \alpha_2 \text{college}_i +
         \alpha_3 \text{exper}_i + \alpha_4 \text{married}_i +
         \alpha_5 \text{kids}_i + \varepsilon_i,
\]
with \texttt{glm(family = binomial(link = "probit"))}.  Estimates appear in
Table~\ref{tab:probit}.

\begin{table}[h]
    \centering
    \caption{Probit model of union-job choice}
    \label{tab:probit}
    \begin{tabular}{lc}
        \toprule
        & Estimate (s.e.) \\
        \midrule
        (Intercept)         & -6.743*** \\
                            & (0.804)   \\
        \texttt{hgc}        & -1.009*** \\
                            & (0.098)   \\
        \texttt{collegeYes} &  0.397    \\
                            & (0.427)   \\
        \texttt{exper}      &  1.849*** \\
                            & (0.156)   \\
        \texttt{marriedYes} &  0.588**  \\
                            & (0.206)   \\
        \texttt{kids}       &  0.799*** \\
                            & (0.202)   \\
        \midrule
        $N$                 & 2229    \\
        \bottomrule
    \end{tabular}\\[2pt]
    \footnotesize{$**\,p<0.01$, $***\,p<0.001$.  GLM warned of fitted
    probabilities very close to $0$ or $1$, suggesting the linear index is
    able to separate the two groups almost perfectly; the magnitudes of the
    coefficients should therefore be interpreted with caution.}
\end{table}

\section*{9.\ Counterfactual: no penalty for being married or a mother}
Following the procedure in \texttt{example.R}, I (i) compute predicted
probabilities of union membership at the estimated parameters, (ii) zero out
the coefficients on \texttt{marriedYes} and \texttt{kids}, and (iii)
recompute the predicted probabilities.  Averaging over the sample:
\begin{center}
    \begin{tabular}{lc}
        \toprule
        Scenario & Mean predicted Pr(union) \\
        \midrule
        Baseline (estimated coefficients)               & 0.2373 \\
        Counterfactual ($\alpha_4 = \alpha_5 = 0$)      & 0.2174 \\
        \midrule
        Difference (counterfactual $-$ baseline)        & $-0.0199$ \\
        \bottomrule
    \end{tabular}
\end{center}
Because the estimated coefficients on \texttt{married} and \texttt{kids} are
\emph{positive}, removing them lowers the average predicted probability of
choosing a union job by about 2 percentage points (from 23.7\% to 21.7\%).
In other words, the model says that under the status quo, married women and
mothers are \emph{more} likely to take union jobs than their otherwise
identical single, childless counterparts.

\paragraph{Is the model realistic?}  Probably not.  Several features of the
estimates are hard to defend.
\begin{enumerate}\itemsep 2pt
    \item \emph{Sign of \texttt{hgc} and \texttt{exper}.}  More schooling
    sharply lowers the probability of taking a union job ($\hat{\alpha}_1 =
    -1.01$), while more experience sharply raises it ($\hat{\alpha}_3 =
    1.85$).  Both magnitudes are implausibly large for a probit and the
    \texttt{glm.fit} warning about fitted probabilities of $0$/$1$ is a
    classic sign of (near) complete separation -- so the precise numbers
    cannot be trusted.
    \item \emph{No price variable.}  We model the choice between a union and
    non-union job purely from demographics and human-capital variables.  An
    economic model of job choice should at minimum include the union and
    non-union wage offers; the omission means the estimated $\hat{\alpha}$'s
    confound \emph{preferences} with \emph{offer probabilities}.
    \item \emph{Selection bias.}  The sample only contains women who are
    employed, so we are estimating preferences conditional on labour-force
    participation, an issue we just demonstrated is non-trivial in
    Section~7.
    \item \emph{Counterfactual interpretation.}  Even taken at face value,
    the policy ``remove the union-job preference for marriage / motherhood''
    only changes the average probability by $\approx 2$ percentage points.
    The estimated effect is therefore small and, given the issues above,
    not a credible basis for actual policy advice.
\end{enumerate}
A more realistic specification would (a) include union and non-union wage
offers (predicted from the Heckman wage regression), (b) jointly model
labour-force participation and union choice, and (c) be estimated on a
larger and more representative sample.

\section*{10--12.\ Repository synchronisation}
The files \texttt{PS12\_Cheng.R}, \texttt{PS12\_Cheng.tex}, and
\texttt{PS12\_Cheng.pdf} were placed in
\verb|~/DScourseS26/ProblemSets/PS12| on OSCER and synced with the GitHub
fork using the standard \texttt{git add / commit / push} workflow.

\end{document}
